\section{Conclusion and Future Work}
  In our research, we recognize the existing limitations in the
    evaluation of synthetic time series, particularly the absence of a
    universally accepted framework. To address this, we adopt the
    evaluation taxonomy proposed by Stenger et al. and expand upon it to
    create a comprehensive benchmark specifically tailored for Synthetic
    Data Generation for Financial Time Series (SDGFTS).
    Our contributions include the development of a consolidated evaluation
    framework that systematically reviews and integrates various
    assessment methods from leading studies in the field. This approach
    not only enhances the rigor of evaluations but also facilitates the
    comparison of different SDGFTS models, ultimately providing a more
    standardized and reliable means of assessing synthetic data quality.

  We believe that Stonk Bench will serve as a valuable resource for
    researchers and practitioners in the field of synthetic financial
    time series generation.
    We hope that our benchmark will facilitate more rigorous and
    standardized evaluations of SDGFTS models, ultimately advancing
    the state of the art in this important area of research.
    For the broader field of synthetic data generation, our work
    highlights the importance of a comprehensive and unifying
    evaluation frameworks that consider both statistical properties
    and practical utility that is proposed by Stenger \textit{et al.}.
  \subsection{Summary of Findings}
  We presented Stonk-Bench, a comprehensive benchmark for evaluating 
  Synthetic Data Generation for Financial Time Series (SDGFTS). This 
  benchmark integrates statistical fidelity, diversity, stylized facts 
  verification, and utility evaluation through (deep) hedging. We evaluated 
  a suite of 5 traditional, 3 deep learning-based, and 1 other SDGFTS 
  models on the S\&P 500 minute-level dataset.
  \subsection{Stonk Bench Shortcomings}\label{subsub:shortcomings}
      While Stonk Bench represents a significant step forward in the
      evaluation of SDGFTS models, we acknowledge several limitations
      in our current work:
    \begin{enumerate}[label=\textbf{[S\arabic*]}]
      \item \textbf{Incompatibility with other models:}
        Stonk Bench has so far been evaluated mainly on classical
        stochastic models (e.g., GBM, OU, GARCH) and a subset of
        deep-learning approaches (GANs, VAEs).
        Other model families — such as
        reinforcement-learning-based generators, normalizing flows,
        autoregressive networks, and hybrid methods — have not yet
        been benchmarked.
      \item \textbf{No hyperparameter tuning:} For consistency and
        comparability across models, StonkBench evaluations were
        performed using default model configurations, without
        performing hyperparameter optimization. This ensures a fair
        baseline but may not reflect the absolute best performance
        achievable by each model.
      \item \textbf{Metric and procedure generalizability:}
        As a consequence, some metrics or evaluation procedures may
        not directly translate to these untested families, which
        limits the framework's generalizability across all synthetic
        financial time-series generation methods.
      \item \textbf{Data Leakage from Public Datasets:} Current models
        can potentially overfit to publicly available datasets
        \textemdash the same datasets that Stonk Bench is using
        \textemdash leading to inflated performance metrics that do
        not generalize well to unseen data. One possible solution is
        to either curate proprietary datasets or implement a grace
        period so the model can be tested on new data that was not
        publicly available during its development.
      \item \textbf{Computational Resource Requirements:} The
        comprehensive nature of Stonk Bench \textemdash particularly
        accommodating any submitted (deep) hedging model and the deep
        hedging utility evaluation \textemdash demands significant
        computational resources. This may limit accessibility for
        researchers with constrained resources.
    \end{enumerate}
  \subsection{Future Research Directions}
    \subsubsection{Remaining progress} In regard to the next portion
    of our project, we aim to expand Stonk Bench in several key areas.
    \begin{enumerate}[label=\textbf{[F\arabic*]}]
      \item \textbf{Multivariate time series predictions.}
        Extend Stonk Bench to evaluate models on multivariate time
        series data, enabling a more comprehensive assessment of their
        performance in realistic scenarios. We wish test SDGFTS models
        on datasets with multiple correlated assets and datasets with
        Bid-Ask spreads and trading volume to evaluate their ability to
        capture inter-asset dependencies and market microstructure
        effects \cite{Takahashi24}.
        \item \textbf{Data granularity evaluation.}
        We wish to test SDGFTS models across multiple data granularities
        (minute-level, daily, and weekly) to evaluate model performance
        on different temporal scales and their ability to capture both
        short-term and long-term dependencies. The most relevant granularities
        that are commonly used in financial analysis and trading strategies
        are 5-minute, 30-minute, and daily data \cite{Wang20, Cont01}.
    \end{enumerate}

    \subsubsection{Beyond the project} We envision Stonk Bench
    evolving into a community-driven benchmark for SDGFTS evaluation.
    As we acknowledge the shortcomings outlined in
    §\ref{subsub:shortcomings}, we propose the following future
    directions:
    \begin{enumerate}[label=\textbf{[F\arabic*]}, resume]
      \item \textbf{Continuous model inclusion.}
      Maintain Stonk Bench as a living benchmark by incorporating
      newly developed SDGFTS techniques (normalizing flows,
      autoregressive generators, advanced diffusion models,
      reinforcement-learning-based generators, and hybrids)
      \cite{Potluru24}.
      \item \textbf{Rigorous metric selection and ablations.}
      As mentioned in Stenger et al., we wish to rigorously conduct
      ablation studies to identify the most informative and robust
      metrics for different model families and stylized facts;
      produce guidance for a compact, interpretable metric set
      \cite{Stenger24}.
      \item \textbf{Categories of datasets} We wish to format Stonk
      Bench to cover a variaty datasets categorized by asset class (
      equities, forex, commodities, cryptocurrencies), 
      granularity (minute, five-minute, daily), and datatypes (OHCL,
      BAV).
      \item \textbf{Public benchmark platform.}
      Launch a website with documentation, code, datasets, and a
      leaderboard for submitted SDGFTS models; provide CI-enabled
      evaluation pipelines for reproducible comparisons.
      \item \textbf{Community engagement.}
      Foster collaboration via contribution guidelines, open issues,
      and reproducible example workflows so researchers can extend
      and validate the benchmark.
    \end{enumerate}

\begin{acks}
  To professor Irene Huang, University of Toronto, for her
  supervision and guidance throughout the project.
\end{acks}



